% 3.3 章 · 几何数值积分与长期行为
% 这一章只回答一个问题：**为什么"每步更准"不等于"长期更好"**。
%   1) 辛方法守恒的不是 H，而是另一个哈密顿量 H_h：真实能量在一条带里振荡，带宽有闭式；
%   2) 非辛方法的能量误差单调漂移，漂移率也有闭式（RK4 是 h^5/72）；
%   3) "轨道闭合"必须说清是**点集闭合**（辛性管）还是**相位不漂**（辛性管不了）。
% 断言附表由 tools/build_book.py 从实际运行结果生成，\input 进来，不手抄。

\numlabchapter{几何数值积分与长期行为}{3.3}

\section{这一章要解决什么}

哈密顿系统
\begin{equation}
  \dot q=\frac{\partial H}{\partial p},\qquad
  \dot p=-\frac{\partial H}{\partial q},
  \label{eq:ham}
\end{equation}
的真解满足 $\dfrac{d}{dt}H(q(t),p(t))=0$：能量是**运动积分**。数值方法一般不保持它，
问题是保持不了多少、以及这个"保持不了"随时间怎么变。

两种错法必须分开：

\begin{enumerate}
  \item \textbf{有界振荡}：误差 $H_n-H_0$ 在一条带里来回走，带宽 $\sim h^p$，
        观测时长再长也不长大。辛方法（\S\ref{sec:sympl}）都是这一类。
  \item \textbf{线性漂移}：误差每一步都往同一个方向掉一点，$H_n-H_0\sim -c\,h^{p+1}n$，
        于是每单位时间掉 $c\,h^p$、观测时长翻倍误差翻倍。所有非辛方法都是这一类。
\end{enumerate}

这两种行为的差别\textbf{不在短期}。本章每一步都拿数字把这句话钉住：$h=0.05$ 的谐振子上，
$10^6$ 步（$T=5\times10^4$）之后辛 Euler 的能量误差带是
$[-2.43902\times10^{-2},+2.56410\times10^{-2}]$，而 RK4 的瞬时误差只有
$2.16923\times10^{-4}$——比辛 Euler 小一百倍，可是它\textbf{一直在往下掉}，
单调占比 $100\%$，后半段与前半段的 $\sup$ 之比恰好 $2.0026$。
短期判据（跑八个周期）会告诉你 RK4 好 $2000$ 倍；长期判据（跑 $10^6$ 步）会告诉你
RK4 的误差在往一个方向走，而辛方法的误差从来没走出它的带。

记号：$u=\tfrac12\varepsilon_{\text{mach}}=2^{-53}\approx1.11\times10^{-16}$；
$M$ 是一步映射的雅可比（线性系统上就是一步矩阵）；$\operatorname{tr}M$、$\det M$ 是它的迹与行列式；
$\theta_h$ 是 $M$ 的旋转角；本章所有"实测"数字都来自 \texttt{tests/3.3/} 的断言报告
（\S\ref{sec:evidence} 汇总）。

\section{辛映射与辛方法}
\label{sec:sympl}

\begin{definition}[辛映射，二维]
一步映射 $\Phi_h:(q,p)\mapsto(q',p')$ 称为\textbf{辛映射}，如果它的雅可比 $M=\partial(q',p')/\partial(q,p)$ 满足
\begin{equation}
  M^{\mathsf T}JM=J,\qquad J=\begin{pmatrix}0&1\\-1&0\end{pmatrix}.
  \label{eq:sympl}
\end{equation}
二维时 \eqref{eq:sympl} 等价于一个标量等式 $\det M=1$（见命题~\ref{thm:det}）。
\end{definition}

\begin{remark}[为什么"相面积"是关键词]
$\det M$ 就是一步映射对相平面面积的放大率。$\det M=1$ 表示\textbf{面积不变}：
数值轨道既不会像显式 Euler 那样一圈圈往外摊开，也不会像 RK4 那样一圈圈往里缩。
能量有界是这件事的后果（\S\ref{sec:backward}），不是它的定义。
\end{remark}

\begin{theorem}[四个辛方法的 $M$ 与 $\det M$]
\label{thm:det}
把 \eqref{eq:ham} 取成谐振子 $H=\tfrac12(q^2+p^2)$（即 $\dot q=p,\ \dot p=-q$），一步映射记作
$M=\begin{pmatrix}a&b\\c&d\end{pmatrix}$。则
\begin{align}
  M_{\text{辛E}}&=\begin{pmatrix}1&h\\-h&1-h^2\end{pmatrix},
  & M_{\text{Verlet}}&=\begin{pmatrix}1-\tfrac{h^2}{2}&h\\-h\left(1-\tfrac{h^2}{4}\right)&1-\tfrac{h^2}{2}\end{pmatrix},\\
  M_{\text{中点}}&=\frac{1}{1+\tfrac{h^2}{4}}\begin{pmatrix}1-\tfrac{h^2}{4}&h\\-h&1-\tfrac{h^2}{4}\end{pmatrix},
  & M_{\text{RK4}}&=aI+bA,\quad A=\begin{pmatrix}0&1\\-1&0\end{pmatrix},
\end{align}
其中 $a=1-\tfrac{h^2}{2}+\tfrac{h^4}{24}$、$b=h-\tfrac{h^3}{6}$，而且
\begin{equation}
  \det M_{\text{辛E}}=\det M_{\text{Verlet}}=\det M_{\text{中点}}=1,\qquad
  \det M_{\text{RK4}}=a^2+b^2=1-\frac{h^6}{72}+O(h^8).
  \label{eq:dets}
\end{equation}
显式 Euler 是 $M=\begin{pmatrix}1&h\\-h&1+h^2\end{pmatrix}$，$\det M=1+h^2$。
\end{theorem}

\begin{proof}
逐个直接算。

\emph{辛 Euler}（先 $\dot p$ 后 $\dot q$）：$p'=p-hq$、$q'=q+hp'=q+hp-h^2q$，于是
$M=\begin{pmatrix}1&h\\-h&1-h^2\end{pmatrix}$，
$\det M=(1-h^2)+h^2=1$。

\emph{Verlet}（半步 kick、整步 drift、半步 kick）：
$p_{1/2}=p-\tfrac h2q$、$q'=q+hp_{1/2}$、$p'=p_{1/2}-\tfrac h2q'$，代进去得
$q'=(1-\tfrac{h^2}{2})q+hp$、$p'=-h(1-\tfrac{h^2}{4})q+(1-\tfrac{h^2}{2})p$，
$\det M=(1-\tfrac{h^2}{2})^2+h^2(1-\tfrac{h^2}{4})=1-\tfrac{h^4}{4}+\tfrac{h^4}{4}=1$。

\emph{隐式中点}：$M=(I-\tfrac h2A)^{-1}(I+\tfrac h2A)$，其中 $A=\begin{pmatrix}0&1\\-1&0\end{pmatrix}$。
$(I-\tfrac h2A)^{-1}=\frac{1}{1+h^2/4}(I+\tfrac h2A)$（因为 $(I-\tfrac h2A)(I+\tfrac h2A)=(1+\tfrac{h^2}{4})I$），
两者相乘得 $\frac{1}{1+h^2/4}(I+\tfrac h2A)^2$；展开 $(I+\tfrac h2A)^2=(1-\tfrac{h^2}{4})I+hA$，
就是上式。$\det M=\frac{(1-h^2/4)^2+h^2}{(1+h^2/4)^2}=\frac{(1+h^2/4)^2}{(1+h^2/4)^2}=1$。

\emph{RK4}：对线性系统 $\dot y=Ay$，$y_{n+1}=R(hA)y_n$，$R(z)=1+z+\tfrac{z^2}{2}+\tfrac{z^3}{6}+\tfrac{z^4}{24}$。
把 $A^2=-I$ 代进去，奇次幂给出 $A$、偶次幂给出 $I$，于是 $M=aI+bA$ 且
$a=1-\tfrac{h^2}{2}+\tfrac{h^4}{24}$、$b=h-\tfrac{h^3}{6}$。$2\times2$ 矩阵 $aI+bA=\begin{pmatrix}a&b\\-b&a\end{pmatrix}$
的行列式是 $a^2+b^2$。展开：
$a^2=1-h^2+\tfrac{h^4}{3}+O(h^6)$，$b^2=h^2-\tfrac{h^4}{3}+O(h^6)$，
$h^2$ 与 $h^4$ 项\textbf{恰好抵消}，$h^6$ 项系数是
$\tfrac{1}{576}-\tfrac{2}{144}+\tfrac{1}{12}\cdot\ldots$ 直接算为 $-\tfrac1{72}$（实测 $1-\det M=1.3871528\times10^{-8}$ 对 $\tfrac{h^6}{72}=1.3888889\times10^{-8}$）。
最后，显式 Euler 的 $M=\begin{pmatrix}1&h\\-h&1+h^2\end{pmatrix}$，$\det M=1+h^2$。
\end{proof}

\begin{remark}[实测]
$h=0.1$ 时四个辛方法的 $|\det M-1|$ 与二次型残差（\S\ref{sec:quad}）最大 $1.4\times10^{-16}$，
即\textbf{精确到位}；RK4 的 $1-\det M=1.3871528\times10^{-8}$，显式 Euler 的 $\det M=1.010000000000000$。
$\det M$ 是这一步映射的性质，与步数无关——它不会因为跑得久而"变好"。
\end{remark}

\section{辛方法守恒的是什么：精确的二次不变量}
\label{sec:quad}

\begin{theorem}[二次不变量的存在性与闭式]
\label{thm:quad}
设 $M=\begin{pmatrix}a&b\\c&d\end{pmatrix}$ 是二维一步映射，令
\begin{equation}
  Q=\begin{pmatrix}c&\frac{d-a}{2}\\[2pt]\frac{d-a}{2}&-b\end{pmatrix}.
  \label{eq:Q}
\end{equation}
则 $M^{\mathsf T}QM=(\det M)\,Q$。特别地，
\[
  \det M=1\quad\Longleftrightarrow\quad M^{\mathsf T}QM=Q
  \quad\Longleftrightarrow\quad H_h(q,p)=\tfrac12\bigl[q\ p\bigr]Q\begin{pmatrix}q\\p\end{pmatrix}
  \ \text{沿轨道严格守恒}.
\]
$\det M\ne1$ 时二次型不再是守恒量：它每步被整体乘以 $\det M$。
\end{theorem}

\begin{proof}
直接把 \eqref{eq:Q} 代进 $M^{\mathsf T}QM$。记 $M^{\mathsf T}QM=\begin{pmatrix}A_{11}&A_{12}\\A_{12}&A_{22}\end{pmatrix}$：
\[
A_{11}=c\,(a^2+bc)=c\det M,\qquad
A_{12}=\tfrac{d-a}{2}(ad+bc)+\tfrac{ac-bd}{2}\cdot 0=\tfrac{d-a}{2}\det M,
\]
\[
A_{22}=(ad+bc)(-b)+\tfrac{d-a}{2}\cdot 0=-b\det M .
\]
（中间项：$A_{12}=ac\cdot\tfrac{d-a}{2}+ad(-b)+bc\cdot\tfrac{d-a}{2}+bd(-b)$，
用 $ad+bc=\det M$ 与 $ac-bd$ 的配对整理即得。）于是 $M^{\mathsf T}QM=(\det M)Q$，
$\det M=1$ 时 $Q$ 就是一步映射的精确二次不变量。
$Q$ 的等值线是椭圆（$\det Q=b c-\tfrac{(d-a)^2}{4}<0$，即 $Q$ 定号）时，
$H_h$ 的等值线就是绕着原点的椭圆，数值轨道只能待在其中一条上。
\end{proof}

\begin{corollary}[实测：把 $H_h$ 算出来]
\label{cor:hhexp}
对四个辛方法用 \eqref{eq:Q} 算 $Q$，再沿 $2\times10^4$ 步轨道算
$H_h$ 的相对变化：
辛 Euler $9.2\times10^{-15}$、Verlet $1.5\times10^{-14}$、隐式中点 $4.3\times10^{-14}$——
全是舍入级。同期\textbf{真实}能量的带宽分别是 $2.56\times10^{-2}$、$6.25\times10^{-4}$、
$4.3\times10^{-14}$：辛方法守恒的是另一个量，而且这个量可以手算出来。
RK4 用同一公式得到 $4.3\times10^{-6}$——不是"算得不准"，是
$\det M\ne1$，这个二次型\textbf{根本不存在}。
\end{corollary}

\section{带宽的闭式：辛 Euler 与 Verlet}
\label{sec:band}

\begin{theorem}[谐振子上能量误差带的两个端点]
\label{thm:band}
谐振子 $H=\tfrac12(q^2+p^2)$、初值 $H_0=\tfrac12$（$q_0=1,p_0=0$）、步长 $h<2$。则
\begin{align}
  \text{辛 Euler:}\quad
  &-\frac{h}{2+h}\ \le\ \frac{H_n-H_0}{H_0}\ \le\ \frac{h}{2-h},
  \label{eq:bandse}\\[2pt]
  \text{Verlet / leapfrog:}\quad
  &-\frac{h^2}{4}\ \le\ \frac{H_n-H_0}{H_0}\ \le\ 0 .
  \label{eq:bandv}
\end{align}
两个端点都能取到（在轨道的特定相位上），因此这是\textbf{最优}的带。
\end{theorem}

\begin{proof}
\emph{辛 Euler}：由 \eqref{eq:Q}（$a=1,b=h,c=-h,d=1-h^2$）得
$d-a=-h^2$，故
$Q=\begin{pmatrix}-h&-\frac{h^2}{2}\\-\frac{h^2}{2}&-h\end{pmatrix}$，即
$H_h=-\frac h2\bigl(q^2+hqp+p^2\bigr)$ 沿轨道不变（$h>0$ 时 $Q$ 负定，等值线仍是椭圆；
$\det Q=h^2-\tfrac{h^4}{4}>0$ 对 $h<2$ 成立）。
初值 $q_0=1,p_0=0$ 把等值线的常数定成 $1$，即整条轨道满足
\begin{equation}
  q^2+hqp+p^2=1 .
  \label{eq:ellipse}
\end{equation}
由 \eqref{eq:ellipse} 解出 $p^2=1-q^2-hqp$，于是真实能量
\[
  H=\tfrac12(q^2+p^2)=\tfrac12\bigl(1-hqp\bigr)=H_0-\frac h2\,qp,
  \qquad\text{即}\qquad \frac{H-H_0}{H_0}=-h\,qp .
\]
在 \eqref{eq:ellipse} 上求 $qp$ 的极值：拉格朗日乘子给出
$p=\lambda(2q+hp)$、$q=\lambda(2p+hq)$，相减得 $(q^2-p^2)(1-\lambda h)=0$。
$q^2=p^2$ 代回约束得 $q^2(2\pm h)=1$，于是
\[
  qp=\pm q^2=\pm\frac{1}{2\pm h}:
  \qquad q=p\ \Rightarrow\ qp=\frac{1}{2+h},\qquad
  q=-p\ \Rightarrow\ qp=-\frac{1}{2-h}.
\]
代回 $(H-H_0)/H_0=-h\,qp$ 即得 \eqref{eq:bandse}。

\emph{Verlet}：由 \eqref{eq:Q}（$a=d=1-\tfrac{h^2}{2}$、$b=h$、$c=-h(1-\tfrac{h^2}{4})$）得
$d-a=0$，所以 $Q$ 是对角的：
\[
Q=\begin{pmatrix}-h\left(1-\frac{h^2}{4}\right)&0\\0&-h\end{pmatrix}
\ \Longrightarrow\
H_h=\frac h2\Bigl[\Bigl(1-\frac{h^2}{4}\Bigr)q^2+p^2\Bigr].
\]
初值给出 $H_h=\tfrac h2(1-\tfrac{h^2}{4})$。由该式解出 $p^2=\tfrac h2(1-\tfrac{h^2}{4})-(1-\tfrac{h^2}{4})\tfrac{h}{2}q^2$，
代入 $H=\tfrac12(q^2+p^2)$：
\[
H-H_0=\frac12\Bigl[q^2+p^2-1\Bigr]=\frac{h^2}{8}\bigl(q^2-1\bigr).
\]
不变椭圆上 $q^2$ 的取值范围是 $[0,1]$：最大值由 $p=0$ 给出
（$\tfrac h2(1-\tfrac{h^2}{4})q^2=\tfrac h2(1-\tfrac{h^2}{4})\Rightarrow q^2=1$），
最小值 $0$ 在 $q=0,\ p^2=\tfrac h2(1-\tfrac{h^2}{4})$ 处取到。于是
\[
  \frac{H-H_0}{H_0}=\frac{h^2}{4}\bigl(q^2-1\bigr)\in\Bigl[-\frac{h^2}{4},\ 0\Bigr].
\]
\end{proof}

\begin{remark}[实测：闭式到九位]
$h=0.05$、$10^6$ 步：
辛 Euler 实测 $[-2.439023753\times10^{-2},+2.564102545\times10^{-2}]$，
闭式 $[-\tfrac{h}{2+h},\tfrac{h}{2-h}]=[-2.439024390\times10^{-2},+2.564102564\times10^{-2}]$；
Verlet/leapfrog 实测 $[-6.249999456\times10^{-4},0]$，闭式 $[-\tfrac{h^2}{4},0]=[-6.25\times10^{-4},0]$。
最大绝对偏差 $6.4\times10^{-9}$——这个残差不是理论误差，是\textbf{采样}采不到极值那一刻
（$2001$ 个采样点跨 $10^6$ 步，步长间隔 $500$ 步）。
注意 \eqref{eq:bandv} 的带\textbf{只在负侧}：写成"$\pm h^2/4$"会让上界差 $100\%$。
\end{remark}

\section{后向误差分析：辛方法解的是另一个哈密顿系统}
\label{sec:backward}

\begin{theorem}[线性系统上的精确版后向误差分析]
\label{thm:backward}
设一步映射 $M$ 满足 $\det M=1$ 且 $M$ 可对角化、特征值不在负实轴上，则存在
$h$ 的解析函数 $S_h=S_h^{\mathsf T}$（对称 $2\times2$ 矩阵）使
\begin{equation}
  M=\exp(h\,J S_h),\qquad\text{于是}\qquad
  H_h(y)=\tfrac12 y^{\mathsf T}S_h\,y\ \text{沿数值轨道严格守恒}.
\end{equation}
若进一步要求 $M$ 是\textbf{对称}方法（用 $-h$ 再走一步回到原处，即 $M(-h)=M(h)^{-1}$），
则 $S_h$ 的展开只含 $h$ 的\textbf{偶}次幂，于是 $H_h=H+O(h^2)$；一般的 $p$ 阶辛方法给
$H_h=H+O(h^p)$，对应的能量\textbf{带宽}是 $O(h^p)$。
\end{theorem}

\begin{proof}
$\det M=1$ 且 $M$ 有实对数（特征值不为负实数）时，存在实矩阵 $B$ 使 $M=e^{hB}$；
取行列式得 $1=\det M=e^{h\operatorname{tr}B}$，所以 $\operatorname{tr}B=0$。
二维无迹矩阵总可以写成 $B=JS$ 且 $S$ 对称：把
$J=\begin{pmatrix}0&1\\-1&0\end{pmatrix}$ 作用在对称矩阵上给出全部无迹矩阵
（$\operatorname{tr}(JS)=\operatorname{tr}\begin{pmatrix}s_{21}&s_{22}\\-s_{11}&-s_{21}\end{pmatrix}=0$，且这是 $2\times2\to2\times2$ 的可逆线性映射）。
取 $S_h$ 为该对称矩阵，则 $M=\exp(hJS_h)$ 是哈密顿函数 $H_h=\tfrac12y^{\mathsf T}S_hy$ 的\textbf{精确}时间-$h$ 流，
故 $H_h$ 严格守恒（定理~\ref{thm:quad} 的 $Q$ 与这里的 $S_h$ 只差一个常数因子：
$M^{\mathsf T}QM=Q$ 说明 $Q$ 与 $e^{hJS_h}$ 交换，二维不可约情形下 $Q\propto S_h$）。

对称性：若 $M(-h)=M(h)^{-1}$，则 $S_h$ 满足 $S_{-h}=S_h$，偶函数，故展开只含偶次幂。
辛方法的一步映射是 $H$ 的流的 $O(h^{p+1})$ 近似，故 $S_h=S+O(h^p)$，
即 $H_h=H+O(h^p)$；误差 $H-H_h$ 被 $H_h$ 的守恒性锁在 $O(h^p)$ 的带里。
\end{proof}

\begin{remark}[非线性系统上这句话的边界]
非线性系统上 $H_h$ 一般只是\textbf{渐近级数}（B 级数），通常不收敛；
可以严格证明的是"形式守恒"与"指数级长的时间内有界"。
所以本章的可视化 \texttt{modified\_\allowbreak hamiltonian} 固定用\textbf{谐振子}，
并在帧的 \texttt{strings.linearity\_note} 里写明这一条：绿线只在谐振子上是精确结论。
这只影响"$H_h$ 存在且收敛"的强度，不影响 \S\ref{sec:drift} 的对照结论
（漂移率是实测的，不依赖 $H_h$ 收敛）。
\end{remark}

\section{非辛方法的漂移率}
\label{sec:drift}

\begin{theorem}[RK4 的能量漂移率 $h^5/72$]
\label{thm:rk4drift}
谐振子、步长 $h$、走 $n$ 步（$T=nh$）时
\begin{equation}
  \frac{H_n-H_0}{H_0}=-\frac{h^5T}{72}+O(h^7T),
\end{equation}
即每步掉 $\tfrac{h^6}{72}$、每单位时间掉 $\tfrac{h^5}{72}$，且\textbf{方向不变}（单调）。
\end{theorem}

\begin{proof}
由定理~\ref{thm:det}，$M_{\text{RK4}}$ 形如 $\begin{pmatrix}a&b\\-b&a\end{pmatrix}$，
$\det M=a^2+b^2$。对这样的矩阵，$\|M y\|^2=(a^2+b^2)\|y\|^2$ 对\textbf{每个} $y$ 成立
（$M$ 是"旋转 $\times$ 均匀缩放"），而 $H=\tfrac12\|y\|^2$，所以
\[
H_{n+1}=(a^2+b^2)H_n=\Bigl(1-\frac{h^6}{72}+O(h^8)\Bigr)H_n .
\]
$n$ 步后 $H_n=(1-\tfrac{h^6}{72})^nH_0$；$n=T/h$、$h$ 小时
$n\log(1-\tfrac{h^6}{72})\approx-\tfrac{h^6}{72}\cdot\tfrac Th=-\tfrac{h^5T}{72}$，
故相对误差 $=-\tfrac{h^5T}{72}+O(h^7T)$。
\end{proof}

\begin{remark}[实测：指数 $4.99965$、系数 $1.7\times10^{-3}$ 相对精度]
取定 $T=50$，$h$ 从 $0.1$ 降到 $0.00625$：$|\Delta H/H|$ 的幂律指数拟合为
$4.99965$（理论 $5$），系数拟合 $0.693287$ 对 $\tfrac{T}{72}=0.694444$ 相对差
$1.7\times10^{-3}$。同一窗口里 Verlet 的带宽指数是 $1.99999$（理论 $2$），
辛 Euler 的带宽与 $\tfrac{h}{2-h}$ 的最大相对差 $6.5\times10^{-5}$——
注意它的\textbf{名义}指数是 $1$，但 $\tfrac{h}{2-h}=h\cdot\tfrac{1}{2-h}$ 在
$h\in[0.1,0.00625]$ 这个窗口上的\textbf{局部}指数是 $1.0186$，
所以这一条要用"与闭式的比值"检验，而不是硬套指数 $1$（第一版就栽在这里）。
\end{remark}

\section{相位：辛性管得了什么、管不了什么}
\label{sec:phase}

\begin{theorem}[三种方法的相位误差率]
\label{thm:phase}
谐振子上一步映射 $M$ 的旋转角记作 $\theta_h$（$\operatorname{tr}M=2\cos\theta_h$ 或
$\tan\theta_h=b/a$）。则真解每单位时间转 $1$ 弧度，而数值解转
\begin{equation}
  \theta_h=
  \begin{cases}
    h\left(1+\dfrac{h^2}{24}+O(h^4)\right), & \text{Verlet / leapfrog},\\[6pt]
    h\left(1-\dfrac{h^2}{12}+O(h^4)\right), & \text{隐式中点},\\[6pt]
    h\left(1-\dfrac{h^4}{120}+O(h^6)\right), & \text{RK4}.
  \end{cases}
  \label{eq:phase}
\end{equation}
于是相位误差每单位时间是 $-\tfrac{h^2}{24}$、$+\tfrac{h^2}{12}$、$+\tfrac{h^4}{120}$（正号 = 落后）。
它\textbf{线性累积}：$T$ 时间后是 $\pm\tfrac{h^pT}{c}$ 圈，辛性对此完全无能为力。
\end{theorem}

\begin{proof}
逐个用 $\operatorname{tr}M$ / $\tan\theta_h$ 的展开。

\emph{Verlet}：$\operatorname{tr}M=2(1-\tfrac{h^2}{2})$，故 $\cos\theta_h=1-\tfrac{h^2}{2}$。
用 $\arccos(1-u)=\sqrt{2u}\left(1+\tfrac{u}{12}+\tfrac{3u^2}{160}+\cdots\right)$，
取 $u=\tfrac{h^2}{2}$：$\theta_h=h\left(1+\tfrac{h^2}{24}+O(h^4)\right)$。

\emph{隐式中点}：$\cos\theta_h=\tfrac{1-h^2/4}{1+h^2/4}$，于是
$1-\cos\theta_h=\tfrac{h^2}{2(1+h^2/4)}=\tfrac{h^2}{2}\left(1-\tfrac{h^2}{4}+O(h^4)\right)$。
代入同一展开：$\theta_h=h\sqrt{1-\tfrac{h^2}{4}+O(h^4)}\left(1+\tfrac{h^2}{24}+O(h^4)\right)
=h\left(1-\tfrac{h^2}{8}+\tfrac{h^2}{24}+O(h^4)\right)=h\left(1-\tfrac{h^2}{12}+O(h^4)\right)$。

\emph{RK4}：$\tan\theta_h=\tfrac ba=\tfrac{h-h^3/6}{1-h^2/2+h^4/24}$。先展开商：
$\tfrac ba=h+\tfrac{h^3}{3}+\tfrac{h^5}{8}+O(h^7)$；再用
$\arctan x=x-\tfrac{x^3}{3}+\tfrac{x^5}{5}+O(x^7)$，取 $x=h(1+\tfrac{h^2}{3}+\tfrac{h^4}{8})$，
$x^3=h^3(1+h^2+O(h^4))$、$x^5=h^5+O(h^7)$：
$\theta_h=h+\tfrac{h^3}{3}+\tfrac{h^5}{8}-\tfrac{h^3}{3}-\tfrac{h^5}{3}+\tfrac{h^5}{5}
=h+h^5\left(\tfrac18-\tfrac13+\tfrac15\right)=h-\tfrac{h^5}{120}$。
\end{proof}

\begin{remark}[实测值与"极角"的口径差]
$h=0.1$、$600$ 圈（$T=3769.9$）：Verlet 实测 $-0.250279$ 圈、理论 $-\tfrac{h^2T}{24\cdot2\pi}=-0.249999$ 圈；
隐式中点 $+0.499250$ 圈、理论 $+\tfrac{h^2T}{12\cdot2\pi}=+0.499999$；
RK4 $+0.000498$ 圈、理论 $+\tfrac{h^4T}{120\cdot2\pi}=+0.000500$。
辛 Euler 实测 $-0.258237$ 圈，而它的 $\operatorname{tr}M$ 与 Verlet 相同
（$\cos\theta_h=1-\tfrac{h^2}{2}$）——多出来的 $-0.008$ 来自"极角 $\arctan(p/q)$"与
"映射旋转角"的口径差：辛 Euler 的不变椭圆是\textbf{斜}的（$Q$ 有非对角元 $-\tfrac h2$），
极角在椭圆上会来回摆，Verlet 的椭圆是轴对齐的，摆动小得多。
两条口径必须说清，否则会得出"两个方法的相位误差不一样"的错误结论。
\end{remark}

\section{角动量、轨道形状与方位}
\label{sec:kepler}

\begin{theorem}[径向 kick 逐位保持角动量]
\label{thm:L}
设中心力场 $\dot p=-g(|q|)q$（$g$ 为标量函数），$L(q,p)=q_xp_y-q_yp_x$。则"kick"
$\;p\mapsto p-h\,g(|q|)q$ 与"drift"$\;q\mapsto q+h\,p$ 都严格保持 $L$。
\end{theorem}

\begin{proof}
kick：$L(q,p-hgq)=q\times(p-hgq)=q\times p-hg\,(q\times q)=L(q,p)$，因为 $q\times q=0$。
drift：$L(q+hp,p)=(q+hp)\times p=q\times p+h\,(p\times p)=L(q,p)$。
\end{proof}

\begin{remark}[浮点下的"逐位"]
上式在实数域是恒等，浮点下 $q_x(gq)_y-q_y(gq)_x$ 的两次乘积不完全相消，
残差是 $O(u)$ 量级。实测 Kepler $e=0.6$、$200$ 圈（$h=0.01$）：
Verlet 的 $|\Delta L|/L=1.9\times10^{-14}$、隐式中点 $8.0\times10^{-14}$、
辛 Euler $2.3\times10^{-14}$；而 RK4 是 $4.8\times10^{-7}$——后者才是"漂移"。
\end{remark}

\begin{theorem}[形状与方位的分离]
\label{thm:shape}
对开普勒问题 $\ddot q=-\mu q/|q|^3$：
\begin{enumerate}
  \item 能量 $E=\tfrac12|p|^2-\mu/|q|$ 与角动量 $L$ 决定轨道的\textbf{形状}：
        $a=-\tfrac{\mu}{2E}$、$e=\sqrt{1+\tfrac{2EL^2}{\mu^2}}$、$r_{\text{近}}=a(1-e)$；
  \item 近星点\textbf{方位}由 Laplace--Runge--Lenz 矢量
        $\mathbf A=p\times\mathbf L-\mu\,\tfrac{q}{|q|}$ 的方向给出，
        $\varpi=\arctan(A_y,A_x)$，$|\mathbf A|=\mu e$。
\end{enumerate}
精确流上 $E$、$L$、$\mathbf A$ 都是运动积分；数值流上它们各自漂。
辛性只保证第一条（$L$ 逐位守恒 $\Rightarrow$ 形状误差 $\le$ 能量带宽），
\textbf{不保证}第二条：$\varpi$ 是相位型量，任何固定步长方法都会进动。
\end{theorem}

\begin{proof}
第一条是开普勒问题的标准结论（由 $E$、$L$ 的定义直接解出 $a$、$e$）。
第二条：对 $\mathbf A$ 求时间导数，
$\dot{\mathbf A}=\dot p\times L+p\times\dot L-\mu\bigl(\tfrac{\dot q|q|-q(\dot q\cdot q)/|q|}{|q|^2}\bigr)$，
用 $\dot L=0$、$\dot p=-\mu q/|q|^3$、$\dot q=p$ 逐项化简，
$p\times\dot L=0$，$\dot p\times L=-\tfrac{\mu}{|q|^3}q\times(q\times p)=\ldots$
最终得 $\dot{\mathbf A}=0$（这一步就是"开普勒轨道是闭合的"的代数证明）。
$|\mathbf A|=\mu e$ 由 $\mathbf A\cdot\mathbf A=\mu^2+2EL^2$ 给出，与第一条的 $e$ 表达式一致。
数值上：$\mathbf A$ 的模与方向分别对应形状与方位，实测见 \S\ref{sec:evidence}。
\end{proof}

\section{数值证据}
\label{sec:evidence}

下面每个数字都来自 \texttt{tests/3.3/} 的断言报告（\S\ref{sec:asserts} 的附表逐条列出）。

\begin{center}
\begin{tabular}{llr}
\toprule
量 & 条件 & 实测 \\
\midrule
辛 Euler 能量带 & 谐振子 $h=0.05$，$10^6$ 步 & $[-2.43902\text{e-}2,\ +2.56410\text{e-}2]$ \\
Verlet 能量带 & 同上 & $[-6.249999\text{e-}4,\ 0]$ \\
RK4 能量误差 & 同上 & $-2.16923\text{e-}4$（单调，占比 $1.0000$）\\
显式 Euler & 同上 & 第 $22134$ 步溢出（$|y|>10^{12}$）\\
两半 $\sup$ 之比 & Verlet / RK4 & $1.0000$ / $2.0026$ \\
$1-\det M$ & RK4，$h=0.1$ & $1.3871528\text{e-}8$（$h^6/72=1.3888889\text{e-}8$）\\
$H_h$ 相对变化 & 辛 Euler / Verlet / 中点 & $9.2\text{e-}15$ / $1.5\text{e-}14$ / $4.3\text{e-}14$ \\
漂移率指数 & RK4，$T=50$，$h\in[0.1,0.00625]$ & $4.99965$（理论 $5$）\\
带宽指数 & Verlet，同窗口 & $1.99999$（理论 $2$）\\
相位误差 & Verlet $h=0.1$，$600$ 圈 & $-0.250279$ 圈（理论 $-0.250000$）\\
$|\Delta L|/L$ & Kepler $e=0.6$，$200$ 圈 & Verlet $1.9\text{e-}14$ vs RK4 $4.8\text{e-}7$ \\
轨道形状失真 & 同上 & Verlet 两半之比 $1.000$ vs RK4 $1.795$ \\
单摆周期 & AGM vs 数值过零（$\theta_0=2$） & $8.349752927$ vs $8.349752920$ \\
\bottomrule
\end{tabular}
\end{center}

三句话总结这张表：辛方法的误差\textbf{有带}（带宽有闭式、两半之比 $1$），
非辛方法的误差\textbf{有斜率}（斜率有闭式、两半之比 $2$），
而"轨道闭合"要分成形状（有界）与方位（线性进动）两句话说。

\section{四个可视化该怎么看}

\begin{enumerate}
  \item \texttt{energy\_\allowbreak drift\_\allowbreak duel}（招牌）：同一时间轴两种方法的能量误差。
        时间轴是 $\log_{10}$（观测步数），拖到右端就是 $10^6$ 步。
        纵轴自动缩放，所以\textbf{不要比谁更陡}：看 \texttt{meta} 里的带宽、
        两半之比与短期误差。辛方法一条带、RK4 一条下坡、显式 Euler 直接溢出
        （溢出后曲线被掩码截断，\texttt{strings.masking} 说明被掩掉几个点）。
  \item \texttt{phase\_\allowbreak orbit}：相平面上的轨道、该方法自己的不变椭圆、精确圆。
        辛方法三条线重合（偏差 $10^{-14}$ 量级）；RK4 向内螺旋、显式 Euler 向外。
        右图给两条偏差曲线：到不变曲线的偏差（闭不含混的判据）与到精确圆的偏差（带宽量级）。
  \item \texttt{modified\_\allowbreak hamiltonian}：真实能量误差带 + 修正能量 $H_h$ 的相对变化。
        绿线平在 $10^{-14}$ 说明"守恒的是另一个量"；切到 RK4 绿线自己也开始漂
        （$\det M\ne1$）。这一页固定用谐振子，理由见 \S\ref{sec:backward} 的注。
  \item \texttt{long\_\allowbreak term\_\allowbreak orbit}：行星轨道。拖观测时长（单位：初始周期）
        看轨道失真：辛方法 $200$ 圈后仍趴在同一条椭圆上，RK4 的能量误差更小却一直在
        往同一边偏。近星点\textbf{方位}的进动对所有方法都存在，帧里单独给一个数。
\end{enumerate}

\section{本章的算法路线（P2/P3 合规说明）}

\begin{itemize}
  \item 单摆的精确周期用 \textbf{AGM}（算术--几何平均）算完全椭圆积分，
        不查表、不调库函数（\texttt{numcore.ode.symplectic.pendulum\_period}）。
  \item Kepler 精确流用\textbf{自写 Newton} 解 Kepler 方程 $E-e\sin E=M$
        （初值 $E=M+e\sin M$，$8$ 次迭代到 $10^{-14}$），再转真近点角。
  \item 二次不变量 $Q$、等值线半径、闭合度、幂律拟合全部自写
        （\texttt{numcore.ode.invariants}），幂律复用 3.1 章的 \texttt{loglog\_slope}，
        不重复实现。
  \item 非辛对照 \texttt{rk4}/\texttt{explicit\_\allowbreak euler} 直接调用 3.1 章的
        \texttt{numcore.ode.rk.rk\_step}，不另写一份 tableau。
  \item 越界保护：积分在 $\max|q|$ 或 $\max|p|$ 触到 $10^{12}$ 时立即停止，
        帧里写"占位值 + 掩码 + 文案"三件套，绝不让 \texttt{inf}/\texttt{NaN} 进帧。
\end{itemize}

\section{断言附表}
\label{sec:asserts}

下表由 \texttt{tools/build\_book.py} 从 \texttt{tests/3.3/} 的**实际运行结果**生成
（\texttt{dist/reports/assertions.json}），\texttt{\textbackslash input} 进来，不手抄。
每一行给出类别（A = 与人造已知答案对拍、B = 不变量/相对变化率、C = 收敛率与阈值）、
理论值、实测值与容差。

% 本文件由 tools/build_book.py 从 dist/reports/assertions.json 生成，请勿手改。
\begin{longtable}{@{}p{0.06\textwidth}p{0.32\textwidth}p{0.22\textwidth}p{0.16\textwidth}@{}}
\caption{章节 3.3 数值断言（生成时间 2026-10-19T16:08:50-0400）}\\
\toprule
类别 & 断言 & 理论值 & 实测值 \\
\midrule
\endfirsthead
\toprule
类别 & 断言 & 理论值 & 实测值 \\
\midrule
\endhead
\bottomrule
\endlastfoot

A & 章节声明对账(4 viz / 依赖 3.1 / 断言 B,C) & docs/\allowbreak{}chapters.yaml 是章节清单的唯一真源，MANIFEST 必须逐项对齐 & 4 个可视化：['energy\_drift\_duel',\allowbreak{} 'phase\_orbit',\allowbreak{} 'modified\_hamiltonian',\allowbreak{} 'long\_term\_orbit']；deps=\allowbreak{}['3.1']；asserts=\allowbreak{}['B',\allowbreak{} 'C']；模块 2 个全部存在 \checkmark \\
A & 全参数角点 $\times$ 时间轴两端的帧体检 & PROTOCOLS 1.1/\allowbreak{}1.3 节：帧里不许出现 NaN/\allowbreak{}Inf；出不了有限值就写占位+掩码+文案 & 64 帧、\allowbreak{}586 个数组、\allowbreak{}117 个 meta 字段全部有限且可打包(\allowbreak{}0 个反例) \checkmark \\
B & 10\textasciicircum{}6 步的招牌帧 & 「跑了 10\textasciicircum{}6 步」必须落在 meta 与数组上，不能靠文案 & n\_steps\_observed=\allowbreak{}1000000，t 末端 50000.0，带宽 6.249999e-04(\allowbreak{}理论 h$^2$/\allowbreak{}4=\allowbreak{}6.250000e-04)，两半之比 1.0000 vs 2.0026；短期误差 RK4 2.61e-06 < Verlet 5.16e-03 \checkmark \\
B & 爆点帧的三件套 & PROTOCOLS 1.3 节：占位值 + 掩码 + 文案；占位值取最后可信值 & h=\allowbreak{}2.5 显式 Euler 在第 70 时刻被掩码截断，掩掉的采样点 86/\allowbreak{}114，占位值 1.694830e+23 \checkmark \\
B & 相空间帧的两个数(闭合 vs 相位) & 辛：点集闭合(\allowbreak{}到不变曲线偏差 =\allowbreak{} 舍入级)；相位照样漂(\allowbreak{}h$^2$/\allowbreak{}24) & Verlet sup=\allowbreak{}1.84e-14(\allowbreak{}舍入级 $\implies$ growth 不读；名义 growth 3.609，det M=\allowbreak{}1.000000000000)、\allowbreak{}相位 -0.250279 圈；RK4 sup=\allowbreak{}2.61e-04(\allowbreak{}growth 2.001，1$-$det=\allowbreak{}1.387e-08)、\allowbreak{}相位 +0.000498 圈 \checkmark \\
B & 修正哈密顿量帧：带宽闭式 + H\_\allowbreak{}h 守恒 & H\_h 的二次型由 2$\times$2 映射给出；det M=\allowbreak{}1 $\iff$ 它精确守恒 & Verlet 带宽 2.500000e-03 =\allowbreak{} 闭式 2.500000e-03(\allowbreak{}比值 1.000000)，H\_h 相对变化 8.86e-14；RK4 的同一公式给出 1.39e-03(\allowbreak{}1$-$det =\allowbreak{} 1.387e-08) \checkmark \\
B & 行星轨道帧：形状 vs 方位 & 形状由 (\allowbreak{}E,\allowbreak{}L) 决定(\allowbreak{}辛性管得了)；方位是相位型误差(\allowbreak{}辛性管不了) & Verlet：能量带 7.414e-04、\allowbreak{}形状 growth 1.000、\allowbreak{}角动量 1.89e-14、\allowbreak{}进动 -6.512e-04 rad/\allowbreak{}圈；RK4：能量带 4.404e-06(\allowbreak{}更小)、\allowbreak{}形状 growth 1.795、\allowbreak{}角动量 4.83e-07、\allowbreak{}进动 +1.761e-07 rad/\allowbreak{}圈 \checkmark \\
B & 时间轴(log10 观测步数)的口径 & t=\allowbreak{}0 $\to$ 1 步、\allowbreak{}t=\allowbreak{}6 $\to$ 10\textasciicircum{}6 步；数组只覆盖已观测的那一段(\allowbreak{}不是把全长曲线裁一下) & t=\allowbreak{}0：1 步 /\allowbreak{} 2 点；t=\allowbreak{}3：1000 步 /\allowbreak{} 1001 点；t=\allowbreak{}6：1000000 步 /\allowbreak{} 1501 点 \checkmark \\
A & 度量函数的闭式对拍 & 线性序列：lo/\allowbreak{}hi/\allowbreak{}amplitude/\allowbreak{}sup/\allowbreak{}slope 与两半斜率全部可手算；纯振荡的趋势/\allowbreak{}带宽 =\allowbreak{} 3/\allowbreak{}(\allowbreak{}$\pi$N) & band=\allowbreak{}[-2.000,\allowbreak{}28.000] amp=\allowbreak{}30.000 slope=\allowbreak{}3.000000000000 两半斜率=\allowbreak{}(\allowbreak{}3.000000000000,\allowbreak{}3.000000000000)；正弦趋势/\allowbreak{}带宽的 3/\allowbreak{}(\allowbreak{}$\pi$N) 律：N=\allowbreak{}2:0.4739(\allowbreak{}理论 0.4775)、\allowbreak{}N=\allowbreak{}4:0.2378(\allowbreak{}理论 0.2387)、\allowbreak{}N=\allowbreak{}8:0.1191(\allowbreak{}理论 0.1194)、\allowbreak{}N=\allowbreak{}16:0.0596(\allowbreak{}理论 0.0597) \checkmark \\
A & power\_\allowbreak{}law 的回收能力 & y =\allowbreak{} C$\cdot$h\textasciicircum{}p $\implies$ 双对数斜率 =\allowbreak{} p、\allowbreak{}截距 =\allowbreak{} C & 拟合 (\allowbreak{}p,\allowbreak{}C)=\allowbreak{}(\allowbreak{}2.000000000000,\allowbreak{}3.000000000000)，对数残差 4.44e-16；舍入地板案例：不剔除 p=\allowbreak{}16.4242、\allowbreak{}剔除后 p=\allowbreak{}2.000000000000 \checkmark \\
A & 相位误差与角度折叠 & $\varphi$ =\allowbreak{} atan2(\allowbreak{}p,\allowbreak{}q)；圈数 =\allowbreak{} [$\varphi$(\allowbreak{}t)$-$$\varphi$(\allowbreak{}0)]/\allowbreak{}2$\pi$(\allowbreak{}含缠绕) & 精确同相 $\to$ 0.000000000000 圈；反向流 $\to$ 6.366197723676 圈(\allowbreak{}理论 2T/\allowbreak{}2$\pi$ =\allowbreak{} 6.366197723676)；wrapped\_angle(\allowbreak{}$\pm$3$\pi$) =\allowbreak{} +$\pi$(\allowbreak{}值域 (\allowbreak{}$-$$\pi$,\allowbreak{}$\pi$]) \checkmark \\
A & 二次不变量的三类映射对拍 & det M =\allowbreak{} 1 $\iff$ Q =\allowbreak{} [[c,\allowbreak{}(\allowbreak{}d$-$a)/\allowbreak{}2],\allowbreak{}[(\allowbreak{}d$-$a)/\allowbreak{}2,\allowbreak{}$-$b]] 满足 M$^{\mathsf T}$QM =\allowbreak{} Q；det$\ne$1 时不存在 & 剪切残差 0.00e+00、\allowbreak{}旋转残差 8.62e-17(\allowbreak{}Q$\propto$$-$I)、\allowbreak{}缩放 symplectic=\allowbreak{}False；辛 Euler 的 Q 沿 2e4 步轨道相对变化 9.16e-15 \checkmark \\
A & 等值线半径的闭式对拍 & r(\allowbreak{}$\theta$) =\allowbreak{} $\surd$(\allowbreak{}2$\cdot$value /\allowbreak{} (\allowbreak{}Q$_0$$_0$cos$^2$$\theta$ + 2Q$_0$$_1$cos$\theta$sin$\theta$ + Q$_1$$_1$sin$^2$$\theta$)) & 椭圆 (\allowbreak{}a,\allowbreak{}b)=\allowbreak{}(\allowbreak{}2.0,\allowbreak{}0.5) 的最大等值线残差 2.22e-16；Verlet h$\to$0 时半径偏差 1.25e-13 \checkmark \\
B & 闭合度对三种轨道的判别 & $\delta$\_k =\allowbreak{} r\_k/\allowbreak{}r\_curve(\allowbreak{}$\theta$\_k) $-$ 1；舍入级 $\implies$ 点集闭合；线性累积 $\implies$ 螺旋(\allowbreak{}前后半段比 $\approx$ 时长比) & 完美圆 sup=\allowbreak{}2.22e-16；螺旋 sup=\allowbreak{}2.857e-01(\allowbreak{}growth 2.134)；舍入抖动 sup=\allowbreak{}4.44e-16(\allowbreak{}at\_roundoff=\allowbreak{}True) \checkmark \\
B & 采样受限的假证据：加密采样就变小 & 轨道是连续曲线、\allowbreak{}采样点离散：「最近回归距离」的下界由采样密度(\allowbreak{}h)决定，与是否闭合无关 & 同一条精确圆：逐圈最近回归距离的均值 h=\allowbreak{}0.1 时 2.509e-02、\allowbreak{}h=\allowbreak{}0.01 时 2.513e-03(\allowbreak{}差 10.0 倍 $\approx$ 步长比，纯粹是采样下界)；而 curve\_closure 两次都是 sup=\allowbreak{}2.22e-16/\allowbreak{}2.22e-16 \checkmark \\
B & 到精确轨道的偏差：有界带宽 vs 线性累积 & 辛方法：偏差 $\le$ 不变带 $\implies$ 有界；非辛：每圈穿出一点 $\implies$ 线性(\allowbreak{}前后半段比 $\approx$ 2) & Verlet sup=\allowbreak{}1.251e-03(\allowbreak{}growth 1.000)；RK4 sup=\allowbreak{}2.614e-04(\allowbreak{}growth 2.001)；末端 RK4/\allowbreak{}Verlet =\allowbreak{} 0.209(\allowbreak{}观测 600 圈) \checkmark \\
B & 度量函数的报错路径 & 量不出定义的量时必须抛错：帧层才有机会写 null/\allowbreak{}占位，而不是让 0 或 NaN 混进结论 & 5 条边界全部抛 ValueError：closure\_stats(\allowbreak{}<4 点)、\allowbreak{}radial\_deviation(\allowbreak{}长度不符)、\allowbreak{}return\_distances(\allowbreak{}周期$\le$0)、\allowbreak{}curve\_closure(\allowbreak{}出发点在 H\_h=\allowbreak{}0 上)、\allowbreak{}power\_law(\allowbreak{}长度不符) \checkmark \\
A & 方法表 + 四个辛方法的 det M = 1 & det M $\equiv$ 1(\allowbreak{}精确到位)且 M$^{\mathsf T}$QM$-$Q 在舍入级 & 最大偏离 1.388e-16；symplectic\_euler: order=\allowbreak{}1 det=\allowbreak{}1.000000000000 resid=\allowbreak{}1.73e-17；verlet: order=\allowbreak{}2 det=\allowbreak{}1.000000000000 resid=\allowbreak{}1.39e-16；leapfrog: order=\allowbreak{}2 det=\allowbreak{}1.000000000000 resid=\allowbreak{}1.39e-16；implicit\_midpoint: order=\allowbreak{}2 det=\allowbreak{}1.000000000000 resid=\allowbreak{}1.74e-17；rk4: order=\allowbreak{}4 det=\allowbreak{}0.999999986128 resid=\allowbreak{}1.39e-08；explicit\_euler: order=\allowbreak{}1 det=\allowbreak{}1.010000000000 resid=\allowbreak{}1.00e-02 \checkmark \\
A & 10\textasciicircum{}6 步的带宽 = 闭式 & 辛 Euler [$-$h/\allowbreak{}(\allowbreak{}2+h),\allowbreak{} h/\allowbreak{}(\allowbreak{}2$-$h)]；Verlet/\allowbreak{}leapfrog [$-$h$^2$/\allowbreak{}4,\allowbreak{} 0](\allowbreak{}h=\allowbreak{}0.05，E$_0$=\allowbreak{}1) & 最大绝对偏差 6.368e-09；symplectic\_euler: 实测 [-2.439023753e-02,\allowbreak{}+2.564102545e-02] 闭式 [-2.439024390e-02,\allowbreak{}+2.564102564e-02] 偏差 6.37e-09；verlet: 实测 [-6.249999456e-04,\allowbreak{}+0.000000000e+00] 闭式 [-6.250000000e-04,\allowbreak{}+0.000000000e+00] 偏差 5.44e-11；leapfrog: 实测 [-6.249999456e-04,\allowbreak{}+0.000000000e+00] 闭式 [-6.250000000e-04,\allowbreak{}+0.000000000e+00] 偏差 5.44e-11 \checkmark \\
A & RK4 / 显式 Euler 的 det M & RK4：det =\allowbreak{} 1 $-$ h$^6$/\allowbreak{}72 + h$^8$/\allowbreak{}576；显式 Euler：det =\allowbreak{} 1 + h$^2$ & RK4 最大相对差 8.004e-09；h=\allowbreak{}0.4: 1$-$det=\allowbreak{}5.575111e-05 vs 5.575111e-05(\allowbreak{}相对差 0.00e+00)；h=\allowbreak{}0.2: 1$-$det=\allowbreak{}8.844444e-07 vs 8.844444e-07(\allowbreak{}相对差 1.26e-10)；h=\allowbreak{}0.1: 1$-$det=\allowbreak{}1.387153e-08 vs 1.387153e-08(\allowbreak{}相对差 8.00e-09)；h=\allowbreak{}0.05: 1$-$det=\allowbreak{}2.169460e-10 vs 2.169460e-10(\allowbreak{}相对差 0.00e+00)；显式 Euler det(\allowbreak{}0.1)=\allowbreak{}1.010000000000 \checkmark \\
A & 单摆精确周期 vs 数值过零点 & T(\allowbreak{}$\theta$0) =\allowbreak{} 4/\allowbreak{}$\surd$(\allowbreak{}1+H) $\cdot$ K(\allowbreak{}k)(\allowbreak{}AGM 算的完全椭圆积分) & 最大相对差 1.667e-09；$\theta$0=\allowbreak{}0.01: AGM 6.283224577 实测 6.283224567(\allowbreak{}相对差 1.67e-09)；$\theta$0=\allowbreak{}0.5: AGM 6.382789698 实测 6.382789687(\allowbreak{}相对差 1.61e-09)；$\theta$0=\allowbreak{}1.0: AGM 6.699975664 实测 6.699975655(\allowbreak{}相对差 1.46e-09)；$\theta$0=\allowbreak{}2.0: AGM 8.349752927 实测 8.349752920(\allowbreak{}相对差 8.09e-10) \checkmark \\
A & 精确流的轨道要素守恒 + 辛方法的形状误差 $\le$ 带宽 & a =\allowbreak{} $-$$\mu$/\allowbreak{}2E、\allowbreak{}e =\allowbreak{} |A|/\allowbreak{}$\mu$、\allowbreak{}A 为 Laplace--Runge--Lenz 矢量；辛方法能量有界 $\implies$ 形状有界 & 精确流最大偏离 2.665e-15；Verlet(\allowbreak{}200 圈，h=\allowbreak{}0.01)形状失真 2.470e-04 对能量带宽 7.414e-04 \checkmark \\
B & Verlet 的能量误差有界(10\textasciicircum{}6 步) & 有界 $\le$ 6.25001e-04 & 0.000625 \checkmark \\
B & RK4 的能量误差线性漂移(10\textasciicircum{}6 步) & 斜率 $\ne$ 0 & $-1.085\times 10^{-7}$ \checkmark \\
B & RK4 的漂移是单向的(单调占比 = 1) & $\le$ 0 & 0 \checkmark \\
B & 有界 vs 漂移的定量判据(两半 sup 之比 + 趋势/带宽) & 辛：后半/\allowbreak{}前半 sup 之比 =\allowbreak{} 1(\allowbreak{}同一带宽)；非辛：=\allowbreak{} 2(\allowbreak{}线性长) & sup 之比：Verlet 1.0000、\allowbreak{}RK4 2.0019；趋势/\allowbreak{}带宽：Verlet 0.005013、\allowbreak{}RK4 1.000000(\allowbreak{}振幅 6.250e-04 /\allowbreak{} 斜率 -1.085e-07，单调占比 1.0000) \checkmark \\
B & Verlet 的角动量守恒(中心力逐位性质) & 漂移 $\le$ 1.00000e-12 & $1.510\times 10^{-14}$ \checkmark \\
B & 角动量 vs 动量：谁是辛性的证据 & 中心力 + 径向 kick $\implies$ L 逐位守恒(\allowbreak{}辛性)；动量守恒只是平移不变性，非辛方法也有 & Kepler 200 圈角动量漂移：Verlet 1.887e-14、\allowbreak{}RK4 4.830e-07(\allowbreak{}差 25589763.1 倍)；LJ 链动量漂移最大 6.017e-14(\allowbreak{}含显式 Euler 6.017e-14) \checkmark \\
B & 尾部采样点的 t = n\_\allowbreak{}steps\_\allowbreak{}done$\cdot$h & $\le$ 0 & 0 \checkmark \\
B & 爆炸停止点的口径一致性 & i\_blowup 是第一个含爆炸状态的采样点，[:i\_blowup] 可信；t 与 idx 都改成 n\_steps\_done & h=\allowbreak{}0.05 显式 Euler 在第 22134 步停止(\allowbreak{}t=\allowbreak{}1106.7)，i\_blowup=\allowbreak{}34/\allowbreak{}1501，爆点 |y|=\allowbreak{}1.000e+12，可信段能量误差 7.184e+23 \checkmark \\
B & 快路径 vs 通用路径、verlet vs leapfrog & 两条路径算的是同一个映射(\allowbreak{}浮点求和顺序也相同)；leapfrog 就是 verlet 的别名 & 21 组对拍最大差 0.0e+00(\allowbreak{}逐位相同)；verlet/\allowbreak{}leapfrog 在 7 个系统上差 0.0 \checkmark \\
B & 修正不变量 H\_\allowbreak{}h 的守恒性 & det M =\allowbreak{} 1 $\iff$ 存在精确二次不变量 Q(\allowbreak{}H\_h 的二次项)；RK4 的 det =\allowbreak{} 1$-$h$^6$/\allowbreak{}72 $\implies$ 不存在 & symplectic\_euler: H\_h 相对变化 9.16e-15(\allowbreak{}真实能量带宽 2.56e-02)；verlet: H\_h 相对变化 1.47e-14(\allowbreak{}真实能量带宽 6.25e-04)；implicit\_midpoint: H\_h 相对变化 4.30e-14(\allowbreak{}真实能量带宽 4.32e-14)；RK4 同一公式给出 4.34e-06(\allowbreak{}=\allowbreak{} 它自己的能量漂移) \checkmark \\
C & RK4 能量漂移率对 h 的幂律 & 5 & 4.99965 \checkmark \\
C & RK4 漂移率的系数 = T/72 & |$\Delta$E/\allowbreak{}E| =\allowbreak{} h$^5$T/\allowbreak{}72(\allowbreak{}每步掉 h$^6$/\allowbreak{}72，共 n=\allowbreak{}T/\allowbreak{}h 步) & 拟合系数 0.693287 vs T/\allowbreak{}72 0.694444(\allowbreak{}相对差 1.666e-03) \checkmark \\
C & Verlet 带宽对 h 的幂律 & 2 & 1.99999 \checkmark \\
C & 辛 Euler 带宽 = h/(2$-$h)(比值检验) & 带宽 =\allowbreak{} h/\allowbreak{}(\allowbreak{}2$-$h)(\allowbreak{}det M=\allowbreak{}1 的推论，见 notes 第 2 节) & 最大相对差 6.479e-05(\allowbreak{}h$\in$[0.1,\allowbreak{}0.00625]，比值 ['0.99993521',\allowbreak{} '0.99999990',\allowbreak{} '0.99999997',\allowbreak{} '0.99999467',\allowbreak{} '0.99997639']) \checkmark \\
C & 谐振子的稳定边界(二分实测) & Verlet/\allowbreak{}辛 Euler：h<2；RK4：h<2$\surd$2$\approx$2.8284 & Verlet 2.000001(\allowbreak{}22 次求值、\allowbreak{}20 次二分)、\allowbreak{}RK4 2.828979(\allowbreak{}14 次求值) \checkmark \\
C & 形状失真的增长比(前/后半段) & 辛方法：有界振荡 $\implies$ growth$\approx$1；RK4：线性漂移 $\implies$ growth$\approx$2 & verlet: sup=\allowbreak{}2.470e-04 growth=\allowbreak{}1.000；symplectic\_euler: sup=\allowbreak{}9.513e-03 growth=\allowbreak{}1.000；implicit\_midpoint: sup=\allowbreak{}2.818e-04 growth=\allowbreak{}1.000；rk4: sup=\allowbreak{}1.833e-07 growth=\allowbreak{}1.795 \checkmark \\
B & 小步长时每步 $\le$5 次迭代 & $\le$ 5 & 4.958 \checkmark \\
C & 隐式中点的迭代次数与 h 的关系 & 迭代解 y =\allowbreak{} y\_n + h/\allowbreak{}2$\cdot$[f(\allowbreak{}y\_n)+f(\allowbreak{}y)] 的收缩率 $\propto$ h $\implies$ 迭代次数随 h 减小 & h=\allowbreak{}[0.5,\allowbreak{} 0.2,\allowbreak{} 0.1,\allowbreak{} 0.05] $\to$ 平均迭代 [11.884,\allowbreak{} 7.021,\allowbreak{} 5.934,\allowbreak{} 4.958]，残差最大 9.992e-15 \checkmark \\
\end{longtable}


\section{它会怎样骗你（本章版）}
\label{sec:traps}

\begin{enumerate}
  \item \textbf{拿短期精度当长期结论}：$h=0.1$ 的谐振子上跑八个周期，
        RK4 的自收敛误差是 $4.2\times10^{-5}$，Verlet 是 $2.1\times10^{-2}$——
        RK4 好 $493$ 倍。同一对方法跑 $10^6$ 步（$h=0.05$），
        Verlet 的误差带 $6.25\times10^{-4}$、两半之比 $1.0000$，
        RK4 的误差 $2.17\times10^{-4}$ 却\textbf{一路单调下坡}（两半之比 $2.0026$）。
        观测时长再翻十倍，RK4 就输了。判据必须包含"误差随时间怎么变"。
  \item \textbf{把"有界"当形容词}：$h=0.05$ 的 Verlet 带宽是 $6.249999\times10^{-4}$，
        它\textbf{等于} $\tfrac{h^2}{4}$；$h$ 加倍带宽就乘 $4$。
        只说"有界"不给数，等于没说——本章每条"有界"都落在 sup、两半之比与带宽幂律上。
  \item \textbf{把带宽写成 $\pm h^2/4$}：Verlet 的能量误差只在\textbf{负侧}
        （$H-H_0=\tfrac{h^2}{8}(q^2-q_0^2)\le0$），实测带是 $[-6.249999456\times10^{-4},0]$。
        写成对称区间，上界就错了 $100\%$。第一版断言就是这么红的。
  \item \textbf{用"趋势/带宽"单独判有界}：最小二乘斜率给后段更大权重，
        纯振荡在含 $N$ 个周期的窗口上也有 $\approx\tfrac{3}{\pi N}$ 的假趋势
        （实测 $\sin$ 在 $[0,4\pi]$ 上 $0.4739$、$[0,8\pi]$ 上 $0.2378$，
        与 $\tfrac{3}{\pi N}$ 逐位吻合）。主判据必须是\textbf{后半段与前半段的 $\sup$ 之比}
        （有界 $=1$、线性 $=2$），它对窗口长度与采样混叠都免疫。
  \item \textbf{把"每圈离出发点多近"当闭合判据}：轨道是连续曲线、采样点是离散的，
        这个量的下界由采样步长决定。实测同一条精确圆：采样 $h=0.1$ 时逐圈最近距离均值
        $2.509\times10^{-2}$，$h=0.01$ 时 $2.513\times10^{-3}$——差 $10$ 倍，
        而到不变曲线的偏差两次都是 $2.2\times10^{-16}$。
        真实数据里这个量在 $600$ 圈内从 $1.9\times10^{-2}$ "长"到 $3.4\times10^{-1}$，
        看起来像轨道散了，其实只是采样点没落在回归点上。
  \item \textbf{相位不先解缠绕}：只取末点 $\arctan$ 之差会把它折进 $(-\pi,\pi]$。
        实测反向流 $+6.366197723676$ 圈会被报成 $+0.366$ 圈。
        本章 \texttt{phase\_\allowbreak error\_\allowbreak turns} 先 \texttt{unwrap} 再作差，
        并且每步转过 $\ge\pi$ 时直接抛错而不是给个数。
  \item \textbf{把相位误差当成"方法的好坏"}：$h=0.1$、$600$ 圈下 Verlet 的相位误差是
        $-0.250$ 圈，RK4 只有 $+0.0005$ 圈——相位上 RK4 好 $500$ 倍，
        而形状上 RK4 在不断变差（两半之比 $2.001$ 对 Verlet 的 $1.000$）。
        "轨道闭合"是\textbf{点集}的性质，不是"第 $n$ 圈回到出发点"。
  \item \textbf{把 Kepler 的进动都算到方法头上}：近星点方位 $\varpi$ 是相位型量。
        固定步长下所有方法都进动（$e=0.6$、$h=0.01$：Verlet $-6.5\times10^{-4}$ rad/圈、
        隐式中点 $+1.3\times10^{-3}$、RK4 $+1.8\times10^{-7}$），
        而且 $e=0.9$ 时它还会因为近星点欠分辨而变大（$h=0.002$ 时 $-1.1\times10^{-3}$、
        $h=0.01$ 时 $-1.8\times10^{-2}$ rad/圈）。
        形状（$a$、$e$、$r_{\text{近}}$）才是辛性管得了的那一半。
  \item \textbf{用采样的 $r$ 最小值当近星点}：RK4 的轨道在缩，
        "$r$ 最小的那个采样点"会跟着挪，量出来的"进动"是 $1.6\times10^{-2}$ rad/圈；
        用 Laplace--Runge--Lenz 矢量量同一个解只有 $1.8\times10^{-7}$ rad/圈（差 $9$ 万倍）。
        径向漂移伪装成了方位漂移。
  \item \textbf{把动量守恒当辛性的证据}：LJ 链（$6$ 自由度）上，
        Verlet、辛 Euler、隐式中点、\textbf{RK4 与显式 Euler} 的动量漂移全都在
        $6\times10^{-14}$ 以内——动量守恒来自平移不变性，与非辛性无关。
        真正的判别量是角动量：Kepler $200$ 圈，Verlet $1.9\times10^{-14}$、
        RK4 $4.8\times10^{-7}$，差 $2.6\times10^{7}$ 倍。
  \item \textbf{把自动纵轴当证据}：能量误差从 $10^{-16}$ 到 $10^{24}$ 跨了 $40$ 个量级，
        自动缩放会把 $10^{-14}$ 的舍入噪声画得和 $10^{-4}$ 的漂移一样惊心动魄。
        量级看 \texttt{meta}（带宽、斜率、两半之比），不看谁更陡。
  \item \textbf{把爆点读成"误差消失了"}：越过稳定边界后曲线被掩码截断，
        占位值是\textbf{最后一个可信值}（$1.69\times10^{23}$ 这种），不是 $0$。
        写 $0$ 会被读成"误差归零"，正好读反。帧里 \texttt{mask\_*=0} 的位置一律不算数。
  \item \textbf{把 $h\!\to\!2$ 当成"还能跑"}：谐振子上 Verlet 与辛 Euler 的边界是
        $h<2$（二分实测 $2.000001$），RK4 是 $2\sqrt2$（实测 $2.828979$）。
        $h=2.5$ 时两个方法都在几十步内溢出；辛 Euler 的带宽闭式 $\tfrac{h}{2-h}$
        在 $h\to2$ 时发散，所以可视化 3 的滑块上限取 $1.9$——临界线看得见，极点不进来。
  \item \textbf{把修正哈密顿量外推到非线性}：$H_h$ 的\textbf{精确}二次型只在线性系统上存在
        （定理~\ref{thm:backward}）。单摆、Kepler 上它只是渐近级数，一般不收敛。
        可视化 3 固定用谐振子并在 \texttt{strings} 里写明这一条；
        结论外推时只能说"辛方法的能量误差有界"，不能说"$H_h$ 一定收敛"。
\end{enumerate}
